\documentclass[11pt]{article}
\usepackage[T1]{fontenc}
\usepackage{amsmath,amssymb,hyperref}
\usepackage[margin=1in]{geometry}
\title{Interactads: a provisional terminology note}
\author{Contextual Composition Programme}
\date{15 November 2026}
\begin{document}
\maketitle

\paragraph{Purpose and status.}
\emph{Contextual composition} remains the programme's broad name.
\emph{Interactad} is proposed here for a particular presented categorical
package, not as a replacement for the programme name or a new name for
all relative pseudopromonads. This note pins a provisional usage to the
existing construction; it adds no theorem or acceptance obligation.

\section{The package to be named}
Fix the pure base category $C$ and the ambient category-valued relative
module calculus used in the project's foundations. A module has variance
$C^{\mathrm{op}}\times C\to\mathbf{Cat}$, and tensors are in execution order:
\[
 (X\odot_C Y)(A,D)=\int^{B\in C}_{\mathrm{ps}}X(A,B)\times Y(B,D).
\]
The pseudocoend retains native arrows; it inverts the specified pure-middle
balance arrows, not arbitrary native cells.

An \emph{interactad}, in the proposed usage, comprises valid one-sided
context and content presentations, their induced relative pseudopromonads
$M,N$, and a specified full pseudodistributive law
\[
 M(A,B)=(e_A\downarrow B),\qquad
 N(A,B)=(A\downarrow d_B),\qquad
 \tau:N\odot_C M\Longrightarrow M\odot_C N.
\]
The law includes its four invertible unit/multiplication modifications
and all their coherence equations. The crossing $\tau$ itself need not
be invertible. The chosen presentations and pure interface are part of
the package, not information erased by passing to an underlying equivalence
class.

The retained-cell comma normalization identifies the composite by
\[
 M\odot_C N\simeq H,\qquad H(A,B)=(e_A\downarrow d_B).
\]
Its objects are $(P,Q,f:e_AP\to d_BQ)$. At fixed endpoints a native cell
$(\alpha,\beta)$ satisfies $d_B\beta\,f=f'e_A\alpha$.
Changing endpoints are retained too: over $u:A\to A'$ and $v:B\to B'$,
\[
 \alpha:P\to u^*P',\quad \beta:v_!Q\to Q',\qquad
 d\beta\,d\kappa_{v,Q}\,f=f'e_u(P')e_A\alpha.
\]
Here $e_u$ and $\kappa$ are the presentation's context-reindexing and
content-transport comparison maps. In the realized convention the square
is $J(v)\circ h\Rightarrow h'\circ J(u)$.
Relative multiplication is transported through normalization from
\[
 m_H=(m_M\odot m_N)(M\odot\tau\odot N),
\]
with the ambient associators understood; the pure unit is inherited from
the two one-sided units. Thus neither equality of output laws nor an
object-only composition formula supplies the required structure.

\section{Two qualified usages}
\paragraph{Induced interactad.}
An interactad equipped with an identified derivation of its crossing and
coherences from independently stated one-sided hypotheses. This records
how the structure is obtained; it need not distinguish underlying objects
up to equivalence. O1 supplies such a derivation for its stated class:
finite augmented resource families in a coherent symmetric rig, a finitary
cartesian polynomial occurrence presentation, and valid content data with
the specified genuine relative lifts, occurrence compatibility and
Beck--Chevalley transport. The exact O1 hypotheses remain controlling.
Merely providing a crossing does not discharge that independent-construction
result; not every contextad/contentad pair admits one.

\paragraph{Framed interactad.}
An interactad additionally equipped with the full pure-compatible
ordinary-data external action meant in O2. For its Cartesian interface
setting this includes
\[
 \operatorname{st}_X:H(A,B)\to H(X\times A,X\times B),
 \quad
 J(u\times\mathrm{id}_B)\circ\operatorname{st}_Xh
 \cong\operatorname{st}_Yh\circ J(u\times\mathrm{id}_A),
\]
for every $u:X\to Y$, together with strong pure/sequential comparisons,
action coherence, and naturality on native cells and endpoint squares.
A fixed-frame action, a bijection-only action or a merely lax parameter
comparison is not this full structure. Within O2's canonical O1 class,
the full action exists precisely when every canonical transport
$(!_B)_!:\mathcal L_{\varnothing}\to\mathcal L_B$ is an equivalence.
This is a conditional theorem, not an automatic property of interactads.
Here ``framed'' names this specified extra action; it is not an identification
with the independently established notion of a framed bicategory.
The adjectives \emph{induced} and \emph{framed} can apply together.

\section{Theory, realizations, and names left open}
The categorical package must stay distinct from its particular semantic
realizations and their internal theorems, and from applications or finite
diagnostics. Recursive concurrent languages, capability ownership,
probabilistic kernels, schedulers, robot transactions and Byzantine control
can be important and general examples without becoming the definition of
an interactad. A new realization must exhibit the required structure;
its existence or a successful diagnostic alone does not extend the general
theory.

\emph{Contactad} and \emph{contestad} remain unassigned suggestions, perhaps
for future interference and game-theoretic structures respectively.
No such extension is defined by this note, and no existing artifact is
renamed to either term. A distinct name should follow an explicit package
of data, laws and comparison with the existing structure.

A limited exact-name web search, including searches restricted to arXiv,
on 15 November 2026 found no established mathematical usage of these three
candidate names. This is a search outcome, not proof of absence, priority,
or trademark clearance. The new word asserts no novelty for the ambient
relative pseudopromonad or pseudodistributive-law machinery.

\section{Provisional feedback and concurrency vocabulary}
\paragraph{Recommendation, not a completed definition.}
Reserve \emph{reactad} as a short candidate name for a feedback-capable
extension of the presented interaction architecture. \emph{Reactive
interactad} is the explanatory alternative. The name alone does not include
concurrency: \emph{concurrent reactad} would explicitly name the combined
target. These are working names, not claims that the general structures
have been constructed or that every interactad admits the extension.

The other suggestions remain useful descriptive qualifiers:
\begin{itemize}
\item \emph{Iterative interactad}: emphasizes a specified recursion or
iteration operation and its laws; it need not cover every kind of feedback.
\item \emph{Concurrent interactad}: emphasizes compatible concurrent
composition, without presuming feedback.
\item \emph{Dynamic interactad}: a broad umbrella for a specified package
combining these developments, but less informative than naming its actual
operations.
\end{itemize}
These are not asserted synonyms or an established hierarchy. In particular,
adding one operation does not automatically supply another.

\paragraph{Requirements before a definition is pinned.}
A candidate definition should identify the following mathematical data
and laws, rather than importing a particular language or control model:
\begin{enumerate}
\item \textbf{Underlying presentation and recovery.} Specify the base,
context/content presentations, native cells and endpoint transport, and
how the sequential fragment recovers an interactad in the sense above.
If an extension changes the presentation, supply the comparison instead
of silently broadening the word. A normalization reflection alone is not
an equivalence with the original comma presentation.
\item \textbf{Admissible feedback.} Specify the feedback interfaces,
the domain of the closure operation, any delay, seed, guarding or
solution data, and its action on native cells. State its compatibility
with pure substitution, sequential composition and the interaction law,
with the relevant coherence equations. Recursion, instantaneous feedback
and delayed feedback must not be conflated. Neither existence of a
solution to every loop nor an ordinary total trace is implicit.
\item \textbf{Concurrent composition, when claimed.} Specify which
interfaces can be combined, how shared context and retained content are
accounted for, and the comparison cells and coherence relating concurrent
and sequential composition. State whether these comparisons are directed,
invertible, or restricted. A product of independent systems does not by
itself account for interference, and concurrency does not imply that
effects commute.
\item \textbf{Joint compatibility.} For the combined target, feedback
must interact coherently with concurrent composition on its stated domain,
including retained shared effects. Identify what information a permitted
comparison may forget and prove that the operations respect it. Equality
of visible outputs or fixed-input laws alone is not the required native
compatibility.
\end{enumerate}
This is a design brief for a definition, not an axiom system or an added
contract obligation. A full definition still needs precisely typed
operations and explicit coherence diagrams. \emph{Unrestricted concurrency}
should not be a bare adjective: any theorem must state the class of
concurrent assemblies it covers and the restrictions it removes.

\paragraph{Scope discipline.}
Guarded finite-prefix constructions do not by themselves provide an
infinite execution object, an infinite effect, or unguarded feedback.
Concrete probabilistic, nondeterministic, language, scheduler and control
realizations remain separate from these categorical requirements.
Specific fairness, ownership or message-passing mechanisms do not become
core axioms merely because they motivate a realization. No existing
artifact is renamed a reactad by this note.

\paragraph{Naming check.}
A limited exact-name web search on 15 November 2026, including
arXiv- and nLab-restricted searches for \emph{reactad}, found no established
category-theoretic use of that word or of the four qualified interactad
phrases above. It did find unrelated use of
\href{https://www.skypack.dev/view/reactad}{\texttt{reactad} as a software
package name}. This is not exhaustive clearance or a priority claim.
\emph{Interaction category} is already mathematical terminology; see
\href{https://www.kestrel.edu/people/pavlovic/papers/CLCI.pdf}{Pavlovi\'c,
\emph{Categorical Logic of Concurrency and Interaction. I: Synchronous
Processes}}. The proposed names neither rename that theory nor establish
novelty relative to it.

\section*{Source anchors}
The formulas and distinctions above follow the existing programme sources:
\begin{itemize}
\item \emph{Integrated Relative Computation Acceptance Dossier}, sections
``The relative construction and its retained information'',
``The positive theorem and where coherence is proved'', and
``External framing and representation are different operations''.
\item \emph{General Independent Interaction Criterion} (O1),
\texttt{lem:normalization}, the independent-data hypothesis table,
and \texttt{cor:law}.
\item \emph{External Framing and Representation Comparison} (O2),
\texttt{thm:main}, \texttt{thm:transfer}, and \texttt{sec:exact}.
\end{itemize}
These are source anchors for the proposed vocabulary, not a new audit of
every theorem in those documents.
\end{document}
